This appendix \emph{adds} concrete explanations; nothing in the main text depends on reading it. It covers four notions in this order: ideals and annihilators (\S B.1), exponential-polynomials (\S B.2), the Chinese remainder theorem (\S B.3), and the companion matrix on a numerical example (\S B.4).

\subsection*{B.1 Ideals, annihilators and the gcd, starting from integers}
\paragraph{Integers.} The multiples of $6$ form an \emph{ideal} $(6)=\{\dots,-6,0,6,12,\dots\}$: it is closed under addition and under multiplication by any integer. The ideal generated by $4$ and $6$ is $(4,6)=(2)$, since $2=6-4$ and every combination is even: the ideal ``remembers'' only the $\gcd$.

\paragraph{Operators.} Replace integers by operators and ``multiple of $6$'' by ``kills $u$''. For $u=e^{x}$ every operator $Q(\partial_x)$ with $Q(1)=0$ kills $u$, i.e.\ every multiple of $\partial_x-1$. So
\[
\begin{aligned}
\Ann(e^x)&=(\partial_x-1),\\
\Ann(\sin x)&=(\partial_x^2+1),\\
\Ann(xe^x)&=\big((\partial_x-1)^2\big),\\
\Ann(e^x+e^{2x})&=\big((\partial_x-1)(\partial_x-2)\big).
\end{aligned}
\]
In one variable every annihilator ideal has a single generator (the \emph{minimal annihilator}), exactly as every ideal of integers is $(n)$. In two variables an ideal can need several generators, as $(\tau-\tau_0,\xi-\xi_0)$ does for $u=e^{\tau_0t+\xi_0x}$.

\paragraph{The gcd.} $\Ann(u)=G_uJ_u$ separates the part of the ideal that is a single polynomial multiple ($G_u$, ``infinite part'') from the part $J_u$ that has only finitely many common zeros ($J_u$, ``finite debris''). Example: for $u=f(x+t)$ with $f$ generic (not an exponential-polynomial), $\Ann(u)=(\partial_t-\partial_x)$ by Theorem~\ref{thm:ident}, so $G_u=\tau-\xi$ and $J_u=\ring$; for $u=e^{x+t}$, $\Ann(u)=(\tau-1,\xi-1)$, so $G_u=1$ and $J_u=\Ann(u)$.

\subsection*{B.2 Exponential-polynomials: examples and a test}
\begin{center}\small
\begin{tabular}{@{}L{5.4cm}L{2.3cm}L{6.8cm}@{}}
\toprule
Function & In $\EP$? & Reason (derivatives of $u$ span ...)\\\midrule
$xe^{x}$ & yes & $\spn\{xe^x,e^x\}$, dimension $2$\\
$x^2e^{3t}+\cos x$ & yes & finite sum of polynomial $\times$ exponential ($\cos x=\tfrac12(e^{ix}+e^{-ix})$)\\
$e^{-x^2}$ (Gaussian) & no & $\frac{d^k}{dx^k}e^{-x^2}=(\text{polynomial of degree }k)\,e^{-x^2}$: independent for all $k$, infinite dimension\\
$1/(1+x^2)$ & no & derivatives are rational functions with growing denominators: infinite dimension\\
\bottomrule
\end{tabular}
\end{center}
This is Lemma~\ref{lem:ep} in action: ``$u\in\EP$'' means ``its derivatives span a finite-dimensional space''.

\subsection*{B.3 The Chinese remainder theorem, concretely}
\paragraph{Integers.} The system $x\equiv1\pmod2$, $x\equiv2\pmod3$ has the unique solution $x\equiv5\pmod6$: $\Z/6\cong\Z/2\times\Z/3$. Coprimality of $2$ and $3$ is what makes the pair of remainders determine $x$.

\paragraph{One-variable ODE.} Let $p(\partial)=(\partial-1)(\partial-2)$. Then $\ker p=\ker(\partial-1)\oplus\ker(\partial-2)=\spn\{e^x,e^{2x}\}$, and the Lagrange-type projectors are explicit: if $u=ae^x+be^{2x}$,
\[
\begin{aligned}
u_1&=\frac{\partial-2}{1-2}\,u=\;\;a\,e^{x},\\
u_2&=\frac{\partial-1}{2-1}\,u=\;\;b\,e^{2x}.
\end{aligned}
\]
This is Eq.~\eqref{xeq:lagrange} with $\partial_t$ replaced by $\partial$.

\paragraph{Two variables.} For the wave equation $Q_1=\tau-\xi$, $Q_2=\tau+\xi$ are coprime, but the ideal $(Q_1,Q_2)=(\tau,\xi)$ is not the unit ideal: the two lines meet at the origin. The quotient $\ring/(\tau,\xi)=\C$ has dimension $1$, which is the constant ambiguity $\delta(P)=1$ of the wave equation. Theorem~\ref{thm:factor} says that this finite overlap is the \emph{only} failure of the integer picture.

\subsection*{B.4 The companion matrix on a numerical example}
For $u''-3u'+2u=0$ the companion matrix is $C=\begin{pmatrix}0&1\\-2&3\end{pmatrix}$, with characteristic polynomial $\lambda^2-3\lambda+2$ and eigenvalues $1,2$. The projectors of Eq.~\eqref{eq:spectral} are $\Pi_1=\dfrac{C-2\Id}{1-2}$ and $\Pi_2=\dfrac{C-\Id}{2-1}$, and $e^{xC}=e^x\Pi_1+e^{2x}\Pi_2$.

